\section{Empirical Design}

\subsection{Identification}
The object of interest is the return earned in the window that ends before the
FOMC statement is released, in excess of the risk-free rate. Close-to-close daily
returns include the announcement reaction and therefore measure a different
quantity, so the primary test uses intraday prices. The comparison of interest is
between meetings held in calm and in anxious states, where the state is the VIX
close five trading days before the meeting (lagged one week) and the anxious
threshold is fixed at the development-sample median of 15.07 VIX points. The
threshold is not re-estimated on the holdout. The secondary specification
replaces the indicator with the continuous lagged VIX.

Identification rests on the timing of the window and on the state being
predetermined. Lagging the state by a week limits contamination from
anticipation of the meeting itself. The claim is a conditional association
between pre-announcement returns and the prior volatility state, not a causal
effect of the Committee's actions, and the paper is framed accordingly.

\subsection{Threats to inference}
\paragraph{Episode clustering.} Anxious meetings concentrate in a small number of
volatility episodes, so the effective number of independent observations is far
smaller than the number of meetings. We cluster by calendar year as an
ex-ante proxy for episodes, use CR1 cluster-robust standard errors, and treat
the wild cluster bootstrap-t as primary because the holdout has few clusters.
Asymptotic $t$ statistics are descriptive only.

\paragraph{Power.} With a small number of independent episodes the design can
detect a large premium but cannot distinguish a half-sized premium from zero. We
therefore report minimum detectable effects from a clustered power simulation
using realized volatility, and present non-rejections as bounds on the premium
rather than as evidence of absence.

\paragraph{Release-time measurement.} Statement release times are assigned by
rule rather than observed per meeting, so the window may include a few minutes of
the announcement reaction in some early meetings. This affects the level of the estimate more than the state comparison.
\emph{Update (verified against Federal Reserve release pages).} Eight development meetings
(2011-04-27, 2011-06-22, 2011-11-02, 2012-01-25, 2012-04-25, 2012-06-20, 2012-09-13,
2012-12-12) released the statement at about 12:20--12:35~ET, not 14:15~ET, so their 14:00 window
end follows the announcement. The holdout (2020 onward) is unaffected: its statements are
released at 14:00~ET. Appendix~\ref{app:sens} reports the development-sample effect of this error.

\paragraph{Non-FOMC confounds.} A premium measured on FOMC days could reflect
generic high-volatility or pre-event compensation. The matched-VIX placebo
compares FOMC windows with non-FOMC windows on days with similar lagged VIX, so a
Fed-specific premium requires an excess over matched days.

\paragraph{State measurement.} VIX comparability has changed over time as
short-dated option volume has grown. This matters mainly for the exploratory
VIX-change test, which is outside the frozen specification.

\paragraph{Specification search.} The specification, threshold, outcome window and
inference procedure are fixed in the pre-analysis plan before any estimate on
post-2019 data, and the holdout is run once. Tests outside the plan are labelled
exploratory and cannot change the primary conclusion.
